\documentclass[10pt,a4paper]{amsart}
\usepackage{amsmath,amssymb,amsthm,graphicx}
% Retained minimal preamble for the Stacks Project chapter more-algebra.tex (More on Algebra).
% The chapter file pinned-original-source.tex carries only \input{preamble} on its first line, so no
% preamble text lives inside the chapter itself.  The Stacks Project uses one single shared preamble for
% every chapter; the snapshot stages it identically as sources/stacks-chow/preamble.tex and
% sources/stacks-etale-cohomology/preamble.tex.  The block below is copied mechanically, line for line,
% from sources/stacks-chow/preamble.tex lines 14-18 (the Xy-pic setup) and 154-217 (the native theorem
% declarations together with the mathematical macros).  Nothing was added, removed or redefined; the
% wrapper class and the amsmath/amssymb/amsthm/graphicx base packages are supplied by the assembler.
% Chapter: more-algebra.tex -- More on Algebra.
% For commutative diagrams we use Xy-pic
\usepackage[all]{xy}

% We use 2cell for 2-commutative diagrams.
\xyoption{2cell}

% Theorem environments.
%
\theoremstyle{plain}
\newtheorem{theorem}[subsection]{Theorem}
\newtheorem{proposition}[subsection]{Proposition}
\newtheorem{lemma}[subsection]{Lemma}

\theoremstyle{definition}
\newtheorem{definition}[subsection]{Definition}
\newtheorem{example}[subsection]{Example}
\newtheorem{exercise}[subsection]{Exercise}
\newtheorem{situation}[subsection]{Situation}

\theoremstyle{remark}
\newtheorem{remark}[subsection]{Remark}
\newtheorem{remarks}[subsection]{Remarks}

\numberwithin{equation}{subsection}

% Macros
%
\def\lim{\mathop{\mathrm{lim}}\nolimits}
\def\colim{\mathop{\mathrm{colim}}\nolimits}
\def\Spec{\mathop{\mathrm{Spec}}}
\def\Hom{\mathop{\mathrm{Hom}}\nolimits}
\def\Ext{\mathop{\mathrm{Ext}}\nolimits}
\def\SheafHom{\mathop{\mathcal{H}\!\mathit{om}}\nolimits}
\def\SheafExt{\mathop{\mathcal{E}\!\mathit{xt}}\nolimits}
\def\Sch{\mathit{Sch}}
\def\Mor{\mathop{\mathrm{Mor}}\nolimits}
\def\Ob{\mathop{\mathrm{Ob}}\nolimits}
\def\Sh{\mathop{\mathit{Sh}}\nolimits}
\def\NL{\mathop{N\!L}\nolimits}
\def\CH{\mathop{\mathrm{CH}}\nolimits}
\def\proetale{{pro\text{-}\acute{e}tale}}
\def\etale{{\acute{e}tale}}
\def\QCoh{\mathit{QCoh}}
\def\Ker{\mathop{\mathrm{Ker}}}
\def\Im{\mathop{\mathrm{Im}}}
\def\Coker{\mathop{\mathrm{Coker}}}
\def\Coim{\mathop{\mathrm{Coim}}}

% Boxtimes
%
\DeclareMathSymbol{\boxtimes}{\mathbin}{AMSa}{"02}

%
% Macros for moduli stacks/spaces
%
\def\QCohstack{\mathcal{QC}\!\mathit{oh}}
\def\Cohstack{\mathcal{C}\!\mathit{oh}}
\def\Spacesstack{\mathcal{S}\!\mathit{paces}}
\def\Quotfunctor{\mathrm{Quot}}
\def\Hilbfunctor{\mathrm{Hilb}}
\def\Curvesstack{\mathcal{C}\!\mathit{urves}}
\def\Polarizedstack{\mathcal{P}\!\mathit{olarized}}
\def\Complexesstack{\mathcal{C}\!\mathit{omplexes}}
% \Pic is the operator that assigns to X its picard group, usage \Pic(X)
% \Picardstack_{X/B} denotes the Picard stack of X over B
% \Picardfunctor_{X/B} denotes the Picard functor of X over B
\def\Pic{\mathop{\mathrm{Pic}}\nolimits}
\def\Picardstack{\mathcal{P}\!\mathit{ic}}
\def\Picardfunctor{\mathrm{Pic}}
\def\Deformationcategory{\mathcal{D}\!\mathit{ef}}

% --- appended 2026-10-12 (Chief): the project own definitions of the reference/slogan/history
% display environments, copied mechanically from the shared Stacks preamble lines 8-12
% (\usepackage{verbatim} plus the three \newenvironment lines, which the chapter preamble keeps
% verbatim from the project preamble). Nothing added, removed or redefined: these environments
% carry no mathematical content and the project itself typesets them as comments. Without these
% three lines a counted statement whose own text uses them fails with Environment slogan undefined.
\usepackage{verbatim}
\newenvironment{reference}{\comment}{\endcomment}
\newenvironment{slogan}{\comment}{\endcomment}
\newenvironment{history}{\comment}{\endcomment}
\usepackage{cleveref}
\crefname{originallemma}{original lemma}{original lemmas}
\Crefname{originallemma}{Original lemma}{Original lemmas}
\makeatletter
\newlabel{lemma-K-flat-resolution}{{14380}{source}{}{sourcepointer-1}{}}
\newlabel{lemma-K-flat-resolution@cref}{{[originallemma][0][]14380}{[1][0][]source}{}{}{}}
\makeatother
\crefname{originallemma}{original lemma}{original lemmas}
\Crefname{originallemma}{Original lemma}{Original lemmas}
\makeatletter
\newlabel{lemma-get-ML-system}{{23780}{source}{}{sourcepointer-2}{}}
\newlabel{lemma-get-ML-system@cref}{{[originallemma][0][]23780}{[1][0][]source}{}{}{}}
\makeatother
\title{Every object of the derived category of an inverse system of rings is represented by a system of K-flat complexes}
\author{The Stacks Project Authors}
\date{stacks-project:more-algebra.tex:local-snapshot-20261012; GNU FDL 1.2 or later}
\begin{document}
\maketitle
\setcounter{subsection}{1}
\noindent\textit{Editorial note.} All delivered mathematical text is the author's own, copied from the pinned Stacks Project chapter file more-algebra.tex; no statement or proof step was rewritten, reordered or removed.\par\medskip
\noindent\textit{Editorial note.} The chapter file pinned-original-source.tex carries only \textbackslash{}input\{preamble\} on its first line, so no preamble text lives inside the chapter. The retained file sources/stacks-more-algebra/retained-preamble.tex is a mechanical, line-for-line copy of the project's single shared preamble as staged in the snapshot; nothing was added, removed or redefined and no package was added.\par\medskip
\noindent\textit{Editorial note.} The source URL is the source-level Stacks Project URL, as required when more than one card is produced from one source; the per-claim tag is carried by claim\_key.\par\medskip
\noindent\textit{Editorial note.} The author writes '(details omitted)' inside the counted proof span and has the typo 'K\_\{n +1\}\textasciicircum{}\textbackslash{}bullet' with a missing space before the plus in the last paragraph; both are reproduced unchanged. The two references inside the counted proof span are recorded as non-load-bearing navigation pointers.\par\medskip
\noindent\textit{Allowed standard prerequisites (editorial).}\begin{itemize}
\item Representation of an object of the derived category of an inverse system by a system with surjective transition maps. invoked at proof lines 23824-23826. Is the reduction that makes the fibred-product construction available.
\item Existence of K-flat resolutions of complexes. invoked at proof lines 23827-23829 and 23853-23854. Supplies the K-flat complexes at each stage.
\item Termwise surjectivity and the quasi-isomorphism from a fibred product of complexes to the target. used at proof lines 23846-23852. Is the induction step that produces the next compatible stage.
\end{itemize}\medskip
\ifdefined\hypertarget\hypertarget{sourcepointer-1}{}\fi
\noindent\textit{Source reference.} Original lemma 14380: sources/stacks-more-algebra/pinned-original-source.tex, source lines 14380--14380. Same-chapter navigation pointer; the statement it names is standard for the section and the argument that consumes it stays inside the counted proof span.\par\medskip
\ifdefined\hypertarget\hypertarget{sourcepointer-2}{}\fi
\noindent\textit{Source reference.} Original lemma 23780: sources/stacks-more-algebra/pinned-original-source.tex, source lines 23780--23780. Same-chapter navigation pointer; the statement it names is standard for the section and the argument that consumes it stays inside the counted proof span.\par\medskip
\begin{lemma}
\label{lemma-get-K-flat-system}
Let $(A_n)$ be an inverse system of rings. Every
$K \in D(\textit{Mod}(\mathbf{N}, (A_n)))$
can be represented by a system of complexes $(K_n^\bullet)$
such that each $K_n^\bullet$ is K-flat.
\end{lemma}

\begin{proof}
First use Lemma \ref{lemma-get-ML-system} to represent $K$ by a
system of complexes $(M_n^\bullet)$ such that all the transition maps
$M_{n + 1}^\bullet \to M_n^\bullet$ are surjective.
Next, let $K_1^\bullet \to M_1^\bullet$ be a quasi-isomorphism
with $K_1^\bullet$ a K-flat complex of $A_1$-modules
(Lemma \ref{lemma-K-flat-resolution}).
Suppose we have constructed
$K_n^\bullet \to K_{n - 1}^\bullet \to \ldots \to K_1^\bullet$
and maps of complexes $\psi_e : K_e^\bullet \to M_e^\bullet$ such that
$$
\xymatrix{
K_{e + 1}^\bullet \ar[d] \ar[r] & K_e^\bullet \ar[d] \\
M_{e + 1}^\bullet \ar[r] & M_e^\bullet
}
$$
commutes for all $e < n$. Then we consider the diagram
$$
\xymatrix{
C^\bullet \ar@{..>}[d] \ar@{..>}[r] & K_n^\bullet \ar[d]^{\psi_n} \\
M_{n + 1}^\bullet \ar[r]^{\varphi_n} & M_n^\bullet
}
$$
in $D(A_{n + 1})$. As $M_{n + 1}^\bullet \to M_n^\bullet$ is
termwise surjective, the complex $C^\bullet$ fitting into the left
upper corner with terms
$$
C^p = M_{n + 1}^p \times_{M_n^p} K_n^p
$$
is quasi-isomorphic to $M_{n + 1}^\bullet$ (details omitted).
Choose a quasi-isomorphism $K_{n + 1}^\bullet \to C^\bullet$
with $K_{n +1}^\bullet$ K-flat.
Thus the lemma holds by induction.
\end{proof}

\end{document}
